% English translation of checked_pages/esq_119_058.tex; source: cote 119, site 58, PDF page 59, typescript page 56.
% AI-translated by GPT-6 Astra (OpenAI, ChatGPT Pro mode).
\scanpage{119/58}
\mspage{56}

foundations of topology used to formulate it did not exclude
phenomena of wildness. It goes (I hope) without saying
that the need to develop new foundations for
``geometric'' topology by no means excludes the fact that these phenomena,
like everything under the sun, have their reason for being
and their own beauty. More adequate foundations will not eliminate
these phenomena, but will allow us to place them where they
belong, as ``limiting cases'' of phenomena of ``true''
topology.

($^7$) In fact, to reconstruct this system of spaces
\[ (i_o,\ldots,i_n) \longmapsto X_{i_o,\ldots,i_n} \]
contravariant on Drap(I) (for inclusion of flags), it
suffices to know the $X_i$ (or ``\uline{unfolded strata}'') and the $X_{i,j}$
(or ``\uline{connecting tubes}'') for $i,j \in I$, $i < j$, and the morphisms
$X_{i,j} \longrightarrow X_i$ (which are ``bordering'' inclusions)\tnote{TN-058-01}{``Bordering'' renders the author's quoted ``bordantes''; the corresponding adjective appears again with ``embedding'' on the following page.} and $X_{i,j} \rightarrow X_i$
(which are proper fibrations, whose fibers $F_{i,j}$ are
called ``\uline{connecting fibers}'' for the strata with indices $i$ and $j$).
In the case of a tame multiplicity, however, one must
also know the ``\uline{junction spaces}'' $X_{i,j,k}$ $(i < j < k)$ and their
morphisms into $X_{i,j}$, $X_{j,k}$, and above all $X_{i,k}$, fitting into
the following commutative hexagonal diagram, in which the two squares on
the right are cartesian, the arrows $\hookrightarrow$ are immersions (not
necessarily embeddings here), and the other arrows are
proper fibrations:

\diagram{
\begin{tikzcd}[ampersand replacement=\&]
 \& X_{ik} \arrow[rr, hook] \hand{\arrow[dl]} \& \& X_k \& \\
X_i \& \& X_{ijk} \hand{\arrow[ul, hook]} \arrow[rr, hook] \hand{\arrow[dl]} \& \& X_{jk} \hand{\arrow[ul, hook]} \hand{\arrow[dl]} \\
 \& X_{ij} \hand{\arrow[ul]} \arrow[rr, hook] \& \& X_j \&
\end{tikzcd}
}

(NB. This diagram defines $X_{ijk}$ in terms of $X_{ij}$ and $X_{jk}$ over $X_j$,
but not the arrow $X_{ijk} \longrightarrow X_{ik}$, because $X_{ik} \longrightarrow X_{k\unsure{:}}$\tnote{TN-058-02}{Two faint separated marks follow the final subscript k. The colon-shaped reading retains their position; their physical origin and textual role remain unresolved.} is \uline{not} necessarily
an embedding.)

% NOTE: Isolated short dark trace in the lower-right margin; no determinate character or text anchor is established.
\margin{58}
